% !TeX program = xelatex
\documentclass[UTF8,a4paper,12pt]{ctexart}
\usepackage{geometry}
\geometry{left=2.4cm,right=2.4cm,top=2.5cm,bottom=2.5cm}
\usepackage{amsmath,amssymb,bm}
\allowdisplaybreaks[2]
\IfFileExists{/home/lijs/.local/share/fonts/NotoSerifSC-Regular.otf}{%
  \setCJKmainfont[Path=/home/lijs/.local/share/fonts/,BoldFont=NotoSerifSC-Bold.otf]{NotoSerifSC-Regular.otf}%
  \setCJKsansfont[Path=/home/lijs/.local/share/fonts/,BoldFont=NotoSansSC-Bold.otf]{NotoSansSC-Regular.otf}%
}{}
\IfFileExists{newtxtext.sty}{%
  \usepackage{newtxtext}\usepackage{newtxmath}%
}{%
  \usepackage{mathptmx}%
  \IfFontExistsTF{TeX Gyre Termes}{\setmainfont{TeX Gyre Termes}}{%
    \IfFontExistsTF{NimbusRoman-Regular.otf}{%
      \setmainfont{NimbusRoman-Regular.otf}[%
        BoldFont=NimbusRoman-Bold.otf,
        ItalicFont=NimbusRoman-Italic.otf,
        BoldItalicFont=NimbusRoman-BoldItalic.otf]%
    }{\setmainfont{Latin Modern Roman}}%
  }%
}
\usepackage{xcolor}
\usepackage{fancyhdr}
\pagestyle{fancy}
\fancyhf{}
\fancyfoot[C]{\small\thepage}
\renewcommand{\headrulewidth}{0pt}
\renewcommand{\footrulewidth}{0pt}
\usepackage[hidelinks]{hyperref}
\hypersetup{
  pdftitle={第 12、13 章作业（大学物理 II）},
  pdfauthor={李谨硕},
  pdfsubject={胡其图教材：稳恒磁场、磁介质与附加题}
}
\ctexset{
  section={format={\centering\Large\bfseries}},
  subsection={format={\large\bfseries}}
}
\newcommand{\dd}{\,\mathrm{d}}
\newcommand{\ez}{\hat{\bm e}_z}
\newcommand{\er}{\hat{\bm e}_r}
\newcommand{\ep}{\hat{\bm e}_{\varphi}}
\begin{document}

\begin{center}
  {\LARGE\bfseries 第 12、13 章作业}\\[0.55em]
  {\normalsize 大学物理 II（PHY1252H）}\\[0.95em]
  {\small 李谨硕\quad IEEE 计算机方向\quad 电院 2501\quad 2026 年 10 月 28 日}
\end{center}
\vspace{-0.2em}
{\centering\textcolor{black!30}{\rule{0.9\textwidth}{0.8pt}}\par}
\vspace{0.6em}

\section*{胡其图教材：稳恒磁场}

\subsection*{12.13}

解：长直导线在距其 $d$ 处产生的电场和磁场大小分别为
\[
E=\frac{\lambda}{2\pi\varepsilon_0d},
\qquad B=\frac{\mu_0 I}{2\pi d}.
\]
按图示方向，正电荷所受电场力背离导线，磁场力指向导线。要使粒子沿与导线平行的直线运动，需满足
\[
qE=qvB.
\]
因此
\[
\boxed{v=\frac{\lambda}{\varepsilon_0\mu_0I}
=\frac{\lambda c^2}{I}}
\]
速度与导线电流同向，结果与 $q,d$ 无关。对于有质量粒子，还需 $v<c$，即 $\lambda c<I$，才存在物理上可实现的平衡速度。

\subsection*{12.17}

解：取 $I_1$ 所在导线与纸面的交点为原点，$x$ 轴向右，$y$ 轴向上，$z$ 轴垂直纸面向外。图中导线段位于 $x=b$、$0\leq y\leq a$，电流 $I_2$ 沿 $+y$ 方向。

由毕奥--萨伐尔定律，导线段上任一点处的磁场为
\[
\bm B(y)=\frac{\mu_0 I_1}{2\pi(b^2+y^2)}
\left(-y\hat{\bm e}_x+b\hat{\bm e}_y\right).
\]
电流元受到的安培力为
\[
\dd\bm F=I_2\dd y\,\hat{\bm e}_y\times\bm B
=\frac{\mu_0I_1I_2}{2\pi}\frac{y\dd y}{b^2+y^2}\ez.
\]
沿导线积分得
\[
\boxed{\bm F=\frac{\mu_0I_1I_2}{4\pi}
\ln\!\left(1+\frac{a^2}{b^2}\right)\ez}
\]
方向垂直纸面向外。

\subsection*{12.18}

解：圆柱面旋转形成的面电流密度大小为
\[
K=\sigma\omega R.
\]
忽略细长圆柱面的端部效应，它相当于无限长螺线管，内、外磁场分别为
\[
\bm B_{\mathrm{in}}=\mu_0\sigma R\bm\omega,
\qquad \bm B_{\mathrm{out}}=0.
\]
计算柱面所受磁力时，应扣除面元自身的磁场，作用磁场为两侧磁场的平均值：
\[
\bm B_{\mathrm{act}}=
\frac{\bm B_{\mathrm{in}}+\bm B_{\mathrm{out}}}{2}.
\]
因此单位面积上的磁场力为
\[
\bm f=\bm K\times\bm B_{\mathrm{act}},
\qquad
\boxed{\bm f=\frac{\mu_0\sigma^2\omega^2R^2}{2}\er}.
\]
方向沿径向向外，与电荷及旋转方向的正负无关。

\subsection*{12.20}

解：设载流平面本身在上、下两侧产生的磁场分别为 $\bm B_{s+},\bm B_{s-}$。由右手定则，
\[
\bm B_{s+}=-\frac{\mu_0 i}{2}\hat{\bm e}_y,
\qquad
\bm B_{s-}=\frac{\mu_0 i}{2}\hat{\bm e}_y.
\]
外加匀强磁场在两侧相同，因此总磁场满足
\[
\bm B_2-\bm B_1=\bm B_{s-}-\bm B_{s+}
=\mu_0i\hat{\bm e}_y.
\]
代入 $i=100\ \mathrm{A/m}$、$\bm B_1=6.28\times10^{-5}\ez\ \mathrm{T}$，得
\[
\boxed{\bm B_2=
\left(1.257\times10^{-4}\hat{\bm e}_y
+6.28\times10^{-5}\ez\right)\ \mathrm{T}}
\]
\[
\boxed{B_2\approx1.40\times10^{-4}\ \mathrm{T}}.
\]
方向位于 $yOz$ 平面内，与 $+z$ 轴成约 $63.4^\circ$，偏向 $+y$ 方向。

\subsection*{12.21}

解：用叠加法将两阴影导体补成两个完整圆柱，令其电流密度分别为 $\bm j$ 和 $-\bm j$，其中 $j=I/S$。在中部重叠区域，两虚设电流恰好抵消，恢复题设真空区域。

设 $\bm r_1,\bm r_2$ 分别为从 $O_1,O_2$ 指向该区域任一点的位矢。均匀载流圆柱内部磁场为
\[
\bm B_1=\frac{\mu_0}{2}\bm j\times\bm r_1,
\qquad
\bm B_2=-\frac{\mu_0}{2}\bm j\times\bm r_2.
\]
又因为 $\bm r_1-\bm r_2=\overrightarrow{O_1O_2}$，故
\[
\boxed{\bm B=\frac{\mu_0}{2}\bm j\times\overrightarrow{O_1O_2}},
\qquad
\boxed{B=\frac{\mu_0Id}{2S}}.
\]
真空区域内为匀强磁场，方向垂直于电流方向和两轴线连线，由 $\bm j\times\overrightarrow{O_1O_2}$ 确定，其中 $\bm j$ 取 $O_1$ 一侧导体的电流密度方向。

\subsection*{12.35}

解：取无限长导线为 $y$ 轴，$x$ 轴由导线指向 $AB$。在导线右侧，$I_1$ 产生的磁场垂直纸面向里，大小为
\[
B(x)=\frac{\mu_0I_1}{2\pi x}.
\]
图中 $I_2$ 从 $A$ 流向 $B$，故电流元所受安培力向上，且
\[
\dd F=I_2B(x)\dd x.
\]
于是
\[
F=\frac{\mu_0I_1I_2}{2\pi}\int_d^{d+L}\frac{\dd x}{x},
\qquad
\boxed{F=\frac{\mu_0I_1I_2}{2\pi}\ln\!\frac{d+L}{d}}.
\]
方向平行于长直导线，与图中 $I_1$ 同向。

\subsection*{12.36}

解：取长直导线为 $y$ 轴，圆心为 $(d,0)$，$x$ 轴向右，$z$ 轴垂直纸面向外。按图示，$I_1$ 向下，圆环电流沿逆时针方向。

圆环上各点满足
\[
x=d+R\cos\theta,\qquad y=R\sin\theta,
\qquad B=\frac{\mu_0I_1}{2\pi(d+R\cos\theta)}.
\]
由 $\dd\bm F=I_2\dd\bm l\times\bm B$，得
\[
\dd\bm F=\frac{\mu_0I_1I_2R}{2\pi}
\frac{\cos\theta\hat{\bm e}_x+\sin\theta\hat{\bm e}_y}
{d+R\cos\theta}\dd\theta.
\]
上下对称的电流元的 $y$ 分力抵消。利用
\[
\int_0^{2\pi}\frac{\dd\theta}{d+R\cos\theta}
=\frac{2\pi}{\sqrt{d^2-R^2}}\qquad(d>R),
\]
有
\[
\begin{aligned}
F_x
&=\frac{\mu_0I_1I_2}{2\pi}
\int_0^{2\pi}\left(1-\frac{d}{d+R\cos\theta}\right)\dd\theta\\
&=\mu_0I_1I_2\left(1-\frac{d}{\sqrt{d^2-R^2}}\right).
\end{aligned}
\]
因此
\[
\boxed{F=\mu_0I_1I_2\left(\frac{d}{\sqrt{d^2-R^2}}-1\right)}
\]
合力沿圆心到长直导线的连线，指向长直导线，即相互吸引。

\subsection*{12.37}

解：将圆盘分成半径为 $r$、宽为 $\dd r$ 的同心圆环，其电荷量和等效电流分别为
\[
\dd q=\frac{Q}{\pi R^2}\,2\pi r\dd r,
\qquad
\dd I=\frac{\omega}{2\pi}\dd q.
\]
圆环的磁矩大小为 $\dd m=\pi r^2\dd I$，故
\[
\bm m=\frac{Q R^2}{4}\bm\omega.
\]
于是圆盘受到的磁力矩为
\[
\boxed{\bm\tau=\bm m\times\bm B
=\frac{QR^2}{4}\bm\omega\times\bm B}.
\]
磁场与盘面的夹角为 $\varphi$，所以磁场与盘面法线的夹角为 $\pi/2-\varphi$，力矩大小为
\[
\boxed{\tau=\frac{|Q|\omega BR^2}{4}\cos\varphi}.
\]
方向由 $\bm m\times\bm B$ 确定；按图示取正电荷、$\bm\omega$ 沿 $+y$，则力矩沿 $-z$ 方向。

\subsection*{12.41}

解：取长直导线与线框平面的交点为 $O$，$z$ 轴垂直纸面向外。长直导线的电流 $I_1$ 沿 $+z$，在扇形线框处产生的磁场为
\[
\bm B(r)=\frac{\mu_0 I_1}{2\pi r}\ep.
\]
两段圆弧的电流元与磁场平行或反平行，均不受磁力。上侧径向边（$\theta=\pi/4$）电流向外，下侧径向边（$\theta=-\pi/4$）电流向内，因此
\[
\dd\bm F_+=\frac{\mu_0I_1I_2}{2\pi r}\dd r\,\ez,
\qquad
\dd\bm F_-=-\frac{\mu_0I_1I_2}{2\pi r}\dd r\,\ez.
\]
合力为零。相对于 $O$，两侧力矩分别为
\[
\begin{aligned}
\dd\bm\tau_+&=-\frac{\mu_0I_1I_2}{2\pi}\dd r\,\ep(\pi/4),\\
\dd\bm\tau_-&=\frac{\mu_0I_1I_2}{2\pi}\dd r\,\ep(-\pi/4).
\end{aligned}
\]
由 $\ep(-\pi/4)-\ep(\pi/4)=\sqrt2\hat{\bm e}_x$，得
\[
\boxed{\bm\tau=
\frac{\sqrt2\mu_0I_1I_2}{2\pi}(R_2-R_1)\hat{\bm e}_x}.
\]
力矩沿图中 $x$ 轴正方向。

\section*{胡其图教材：磁介质}

\subsection*{13.2}

\textbf{（1）球形腔表面的磁化电流密度}

解：设 $\er$ 从球心指向腔面，$\theta$ 为 $\bm M$ 与 $\er$ 的夹角。磁介质指向空腔的外法线为 $\bm n=-\er$，由磁化面电流公式
\[
\bm K_m=\bm M\times\bm n,
\]
得
\[
\boxed{\bm K_m=-\bm M\times\er},
\qquad
\boxed{K_m=M\sin\theta}.
\]
若 $\bm M=M\ez$，则 $\bm K_m=-M\sin\theta\,\ep$。

\textbf{（2）磁化电流产生的磁矩}

挖去球形区域等效于在原均匀磁化介质中叠加一个磁化强度为 $-\bm M$ 的球体。因此，腔面磁化电流的磁矩为
\[
\boxed{\bm m=-\frac{4\pi R^3}{3}\bm M}.
\]
其方向与原介质的磁化方向相反。

\subsection*{13.4}

解：图中磁化强度 $\bm M$ 向右，永磁棒左端相当于磁南极，右端相当于磁北极。磁感应线在棒内由左向右，在棒外从右端返回左端。

棒外 $\bm H=\bm B/\mu_0$，两者同向；棒内
\[
\bm B=\mu_0(\bm H+\bm M),
\qquad
\bm H=\frac{\bm B}{\mu_0}-\bm M,
\]
其中退磁场 $\bm H$ 与 $\bm M$ 反向。因此各点方向为
\[
\begin{array}{c|ccccccc}
\text{点} & 1&2&3&4&5&6&7\\ \hline
\bm B & \rightarrow&\leftarrow&\leftarrow&\rightarrow&\rightarrow&\rightarrow&\rightarrow\\
\bm H & \leftarrow&\leftarrow&\leftarrow&\rightarrow&\leftarrow&\leftarrow&\rightarrow
\end{array}
\]
对于细长棒中央的点 1，退磁场很小，常近似取 $\bm H_1\approx0$、$\bm B_1\approx\mu_0\bm M$。题目未给出棒的具体尺寸和各点位置，不能进一步确定各点场强的数值。

\subsection*{13.5}

解：图中点 1 位于永磁环内，点 2 位于狭缝内。设磁环介质部分的平均磁路长度为 $l$，狭缝宽为 $\delta$，忽略漏磁和边缘效应。

磁感应强度的法向分量连续，所以
\[
B_1=B_2=B.
\]
环内和狭缝内分别满足
\[
B=\mu_0(H_1+M),\qquad B=\mu_0H_2.
\]
沿磁路没有自由电流，故
\[
\oint\bm H\cdot\dd\bm l=0,
\qquad H_1l+H_2\delta=0,
\]
其中 $H_1,H_2$ 取沿磁化方向的有符号分量。联立解得
\[
B=\mu_0M\frac{l}{l+\delta},
\qquad
H_1=-M\frac{\delta}{l+\delta},
\qquad
H_2=M\frac{l}{l+\delta}.
\]
对于题设的很狭窄细缝 $\delta\ll l$，场的大小近似为
\[
\boxed{B_1\approx B_2\approx\mu_0M,
\qquad |H_1|\approx0,\qquad |H_2|\approx M}.
\]
$\bm B_1,\bm B_2$ 沿环的磁化方向，缝内 $\bm H_2$ 同向，环内微小的 $\bm H_1$ 反向。

\subsection*{13.7}

解：设螺绕环内总磁感应强度为 $B$。由有磁介质时的安培环路定理，
\[
B=\mu_0\mu_rH,
\qquad B_0=\mu_0H.
\]
磁化电流产生的磁感应强度为
\[
B'=B-B_0=(\mu_r-1)B_0.
\]
代入 $\mu_r=4200$，得
\[
\boxed{\frac{B_0}{B'}=\frac{1}{\mu_r-1}=\frac{1}{4199}},
\qquad \boxed{B_0:B'=1:4199}.
\]

\subsection*{13.8}

\textbf{（1）空间各处的磁感应强度}

解：取中心直导线电流方向为 $+z$，绕轴正方向为 $\ep$。这里 $\mu_1,\mu_2$ 均为绝对磁导率。对半径为 $r$ 的同轴圆周应用安培环路定理：
\[
\oint\bm H\cdot\dd\bm l=I_{\mathrm{free,in}}.
\]
因此
\[
\bm H(r)=
\begin{cases}
\dfrac{I}{2\pi r}\ep,&0<r<R_2,\\[0.5em]
0,&r>R_2.
\end{cases}
\]
由 $\bm B=\mu\bm H$ 得
\[
\boxed{\bm B(r)=
\begin{cases}
\dfrac{\mu_1I}{2\pi r}\ep,&0<r<R_1,\\[0.6em]
\dfrac{\mu_2I}{2\pi r}\ep,&R_1<r<R_2,\\[0.6em]
0,&r>R_2.
\end{cases}}
\]
中心导线按理想线电流处理，$r=0$ 不在上式适用范围内。

\textbf{（2）两磁介质分界面上的磁化电流}

两种介质的磁化强度分别为
\[
\bm M_i=\left(\frac{\mu_i}{\mu_0}-1\right)
\frac{I}{2\pi r}\ep\qquad(i=1,2).
\]
在 $r=R_1$ 处，取从介质 1 指向介质 2 的法线为 $\er$，合磁化面电流密度为
\[
\bm K_m=(\bm M_1-\bm M_2)\times\er
=\frac{(\mu_2-\mu_1)I}{2\pi\mu_0R_1}\ez.
\]
因此该分界面上的总磁化电流为
\[
\boxed{\bm I_{m,R_1}=\frac{\mu_2-\mu_1}{\mu_0}I\,\ez}.
\]
若 $\mu_2>\mu_1$，其方向与中心直导线电流同向；若 $\mu_2<\mu_1$，则反向；若两者相等，则为零。

在外表面 $r=R_2$，还存在介质 2 与真空的分界面，其磁化面电流和总磁化电流分别为
\[
\bm K_{m,R_2}=-\frac{(\mu_2-\mu_0)I}{2\pi\mu_0R_2}\ez,
\qquad
\boxed{\bm I_{m,R_2}=-\left(\frac{\mu_2}{\mu_0}-1\right)I\,\ez}.
\]
对于顺磁质 $\mu_2>\mu_0$，此磁化电流与中心直导线电流反向。

\subsection*{13.9}

\textbf{（1）绕圆盘边缘的磁化电流}

解：圆盘半径 $R=0.040\ \mathrm{m}$、厚度 $h=0.0050\ \mathrm{m}$，磁化强度为 $M=1.4\times10^6\ \mathrm{A/m}$。取磁化方向为 $+z$。

上下端面的法线平行于 $\bm M$，故无磁化面电流；圆盘侧面的磁化面电流密度为
\[
\bm K_m=\bm M\times\er=M\ep.
\]
因此绕圆盘边缘的总磁化电流为
\[
\boxed{I_m=Mh=7.0\times10^3\ \mathrm{A}}.
\]
电流方向与 $\bm M$ 构成右手关系。

\textbf{（2）圆盘中心的磁感应强度}

侧面的磁化电流等效为长度 $h$、半径 $R$ 的有限长螺线管。位于高度 $z$、厚度 $\dd z$ 的圆环电流为 $\dd I=M\dd z$，在圆盘中心产生的磁场为
\[
\dd B=\frac{\mu_0MR^2\dd z}{2(R^2+z^2)^{3/2}}.
\]
所以
\[
\begin{aligned}
B_O
&=\frac{\mu_0MR^2}{2}
\int_{-h/2}^{h/2}\frac{\dd z}{(R^2+z^2)^{3/2}}\\
&=\frac{\mu_0Mh}{2\sqrt{R^2+(h/2)^2}}.
\end{aligned}
\]
代入得
\[
\boxed{\bm B_O\approx0.110\ez\ \mathrm{T}}.
\]
若将薄圆盘边缘电流近似为同一平面内的圆电流，也得到 $B_O\approx\mu_0I_m/(2R)=0.110\ \mathrm{T}$。

\textbf{（3）圆盘中心的磁场强度}

由 $\bm B=\mu_0(\bm H+\bm M)$，有
\[
\bm H_O=\left(\frac{B_O}{\mu_0}-M\right)\ez.
\]
代入计算得
\[
\boxed{\bm H_O\approx-1.31\times10^6\ez\ \mathrm{A/m}}.
\]
磁场强度方向与磁化方向相反，为退磁场。

\section*{附加题}

\textbf{载流圆环的一小段圆弧受到其余部分的磁作用力}

解：将圆环视为理想线电流，设所取圆弧的圆心角为 $\Delta\theta=2\alpha\ll1$，圆弧中点位于 $x$ 轴正方向。圆环位于 $xy$ 平面，电流沿逆时针方向；所取圆弧为 $-\alpha\leq\theta\leq\alpha$，其余部分为 $\alpha\leq\phi\leq2\pi-\alpha$。

\textbf{（1）其余圆弧在所取圆弧上产生的磁场}

源电流元与场点的距离为
\[
s=2R\sin\frac{\phi-\theta}{2}.
\]
由毕奥--萨伐尔定律，其磁场垂直圆环平面向外，且
\[
\dd B_z=\frac{\mu_0I}{16\pi R}
\frac{\dd\phi}{\sin[(\phi-\theta)/2]}.
\]
积分得
\[
\begin{aligned}
B_z(\theta)
&=\frac{\mu_0I}{16\pi R}
\int_\alpha^{2\pi-\alpha}
\frac{\dd\phi}{\sin[(\phi-\theta)/2]}\\
&=\frac{\mu_0I}{8\pi R}
\ln\!\left[\cot\frac{\alpha+\theta}{4}
\cot\frac{\alpha-\theta}{4}\right].
\end{aligned}
\]

\textbf{（2）整段小圆弧所受合力}

电流元的受力为
\[
\dd\bm F=IRB_z(\theta)
\left(\cos\theta\hat{\bm e}_x+\sin\theta\hat{\bm e}_y\right)\dd\theta.
\]
由对称性 $F_y=0$，所以
\[
F_x=\frac{\mu_0I^2}{8\pi}
\int_{-\alpha}^{\alpha}\cos\theta
\ln\!\left[\cot\frac{\alpha+\theta}{4}
\cot\frac{\alpha-\theta}{4}\right]\dd\theta.
\]
利用对称性和分部积分，积分结果为
\[
\int_{-\alpha}^{\alpha}\cos\theta
\ln\!\left[\cot\frac{\alpha+\theta}{4}
\cot\frac{\alpha-\theta}{4}\right]\dd\theta
=4\sin\alpha\left(\ln\cot\frac{\alpha}{2}+1\right).
\]
于是整段圆弧的合力为
\[
\boxed{\bm F=\frac{\mu_0I^2}{2\pi}
\sin\frac{\Delta\theta}{2}
\left(\ln\cot\frac{\Delta\theta}{4}+1\right)
\hat{\bm e}_x}.
\]
对于题设 $\Delta\theta\ll1$，可写为
\[
\boxed{\bm F=\frac{\mu_0I^2\Delta\theta}{4\pi}
\left[\ln\frac{4}{\Delta\theta}+1
+O\!\left((\Delta\theta)^2\ln\frac{1}{\Delta\theta}\right)\right]
\hat{\bm e}_x}.
\]
方向沿圆弧中点的径向向外，即圆环有膨胀趋势。

\textbf{（3）小电流元的常用近似}

若只保留 $\Delta\theta\ln(1/\Delta\theta)$ 的首项，则
\[
\boxed{F\sim\frac{\mu_0I^2\Delta\theta}{4\pi}
\ln\frac{4}{\Delta\theta}}.
\]
直接将小圆弧当作位于其中点的电流元，也给出这一对数主项。其余圆弧在两端附近的磁场并不均匀，因此对整段圆弧积分时还会产生上式中的 $+1$ 修正。

理想线电流模型在端点处的磁场有对数奇性，但上述整段合力积分收敛。若讨论真实导线的局部应力，还需给定导线截面等信息；上述结果适用于导线足够细、所取弧长大于其横向尺寸的模型。

\end{document}
